% ==============================================================================
% SEZIONE 2: FRAMEWORK TEORICI PER L'ANALISI STRATEGICA
% ==============================================================================

\section{Framework Teorici}

\begin{frame}{Framework Teorici per l'Analisi Strategica}
    La navigazione nelle dinamiche del settore utility richiede modelli strutturati capaci di decodificare le interdipendenze tra forze macro e decisioni micro.
    \vspace{0.2cm}
    
    \note{I modelli strutturati decodificano le interdipendenze tra forze macroeconomiche e decisioni microaziendali.}

    \begin{block}{Analisi PESTEL e Rischio Paese}
        Essenziale per la mappatura del rischio paese. Monitorare le divergenze tra PIL (\textit{GDP}) e inflazione (\textit{CPI}) per calibrare l'allocazione del capitale.
    \end{block}
    \vspace{0.1cm}
    \pause
    \note{L'analisi PESTEL mappa il rischio paese valutando le divergenze tra GDP e CPI per allocare il capitale.}

    \begin{block}{Value Chain del Settore Energetico}
        Pivot strategico dalla Generazione Flessibile e dal Trading verso la centralità delle \textbf{Grids (Reti)}. Lo spostamento dell'EBITDA verso asset regolati isola il valore dalle fluttuazioni wholesale.
    \end{block}
    \vspace{0.1cm}
    \pause
    \note{La catena del valore evidenzia un pivot verso le grids. Spostare l'EBITDA su asset regolati protegge dal mercato wholesale.}

    \begin{block}{Framework dell'Innovazione Digitale}
        Il passaggio dal consumatore passivo al \textbf{prosumer} è abilitato dai contatori \textbf{2G}. La granularità del dato (15 min) stabilizza la rete e abilita servizi a valore aggiunto.
    \end{block}
    \note{L'innovazione digitale abilita il passaggio al prosumer tramite contatori 2G, stabilizzando la rete con dati ogni 15 minuti.}
\end{frame}


\begin{frame}{Vantaggi dei Framework Strutturati}
    L'adozione di un approccio metodologico rigoroso garantisce tre vantaggi competitivi nell'analisi delle grandi utility:
    \vspace{0.3cm}

    \begin{itemize}
        \item \textbf{Mappatura del Rischio Geografico:} Identificazione tempestiva delle asimmetrie macroeconomiche tra le diverse aree Tier 1.
        \pause
        \note{I framework strutturati offrono tre vantaggi: primo, la mappatura delle asimmetrie macroeconomiche tra aree Tier 1.}
        
        \item \textbf{Visione Olistica dell'Asset Base:} Integrazione sinergica e strategica tra la capacità produttiva fisica e l'infrastruttura digitale intangibile.
        \pause
        \note{Secondo, l'integrazione sinergica tra capacità produttiva fisica e infrastruttura digitale intangibile.}
        
        \item \textbf{Visibilità dei Margini:} Distinzione netta tra ricavi ad alta volatilità (Trading) e flussi di cassa stabili derivanti dalla \textbf{RAB (Regulatory Asset Base)}.
    \end{itemize}
    \note{Terzo, la visibilità dei margini separando il trading volatile dai flussi di cassa stabili della RAB. Passiamo ai market drivers.}
\end{frame}
